\begin{proof}[Proof of \cref{obj:lem:causal-realization}]
\begin{enumerate}
\item
Throughout, \(j\in\{1,\ldots,d\}\), and all treatment and outcome indices range over \(\{0,1\}\). Recall the cell quantities
\[
\begin{aligned}
p_j(P)&=P(X=j),\\
s_{aj}(P)&=P(X=j,A=a),\\
q_{aj}(P)&=P(X=j,A=a,Y=1),\\
\mu_{aj}(P)&=
\begin{cases}
q_{aj}(P)/s_{aj}(P),&s_{aj}(P)>0,\\
0,&s_{aj}(P)=0.
\end{cases}
\end{aligned}
\]
The propensity convention in \cref{obj:synth_2} is
\[
e_j(P)=
\begin{cases}
s_{1j}(P)/p_j(P),&p_j(P)>0,\\
1/2,&p_j(P)=0.
\end{cases}
\]
Since these quantities are obtained by summing probabilities,
\[
p_j(P)\ge0,\qquad
\sum_{j=1}^{d}p_j(P)=1,\qquad
s_{0j}(P)+s_{1j}(P)=p_j(P),\qquad
0\le q_{aj}(P)\le s_{aj}(P).
\]
Thus \(e_j(P)\in[0,1]\): on an occupied cell this follows from
\(0\le s_{1j}(P)\le p_j(P)\), and on a null cell it follows from the
value \(1/2\). Similarly, division by \(s_{aj}(P)>0\), together with
the stipulated value on a null arm, gives
\[
0\le\mu_{aj}(P)\le1.
\]

Define the treatment and potential-outcome Bernoulli masses by
\[
\beta^A_j(a)=
\begin{cases}
1-e_j(P),&a=0,\\
e_j(P),&a=1,
\end{cases}
\qquad
\beta^Y_{aj}(z)=
\begin{cases}
1-\mu_{aj}(P),&z=0,\\
\mu_{aj}(P),&z=1.
\end{cases}
\]
These definitions apply to every cell and arm, including those with
zero mass. By \cref{obj:synth_7}, they are nonnegative and satisfy
\[
\sum_a\beta^A_j(a)=1,\qquad
\sum_z\beta^Y_{aj}(z)=1.
\]
For a binary pair \((z_0,z_1)\), write
\[
z_a=
\begin{cases}
z_0,&a=0,\\
z_1,&a=1.
\end{cases}
\]
Construct \(H_{\mathrm{ind}}\) through its atom masses:
\[
\Pr_{H_{\mathrm{ind}}}
 (X=j,A=a,Y=y,Y_0=z_0,Y_1=z_1)
=
p_j(P)\beta^A_j(a)
\beta^Y_{0j}(z_0)\beta^Y_{1j}(z_1)
\mathbf1\{y=z_a\}.
\]
Every entry is nonnegative. Summing the factual-outcome indicator
first and then the remaining binary coordinates yields
\[
\begin{aligned}
&\sum_{j=1}^{d}\sum_{a,y,z_0,z_1}
\Pr_{H_{\mathrm{ind}}}
 (X=j,A=a,Y=y,Y_0=z_0,Y_1=z_1)\\
&\quad=
\sum_{j=1}^{d}p_j(P)
 \left(\sum_a\beta^A_j(a)\right)
 \left(\sum_{z_0}\beta^Y_{0j}(z_0)\right)
 \left(\sum_{z_1}\beta^Y_{1j}(z_1)\right)
=1.
\end{aligned}
\]
Hence the displayed table defines a probability law; normalization
preserves its atom masses.
% lean: CausalSmith.Stat.AnnotationRarearmFrontier.independentExtension_fullMass

\item
Summing out both potential outcomes gives
\[
\begin{aligned}
&\Pr_{H_{\mathrm{ind}}}(X=j,A=a,Y=y)\\
&\quad=
\sum_{z_0,z_1}
p_j(P)\beta^A_j(a)
\beta^Y_{0j}(z_0)\beta^Y_{1j}(z_1)
\mathbf1\{y=z_a\}\\
&\quad=p_j(P)\beta^A_j(a)\beta^Y_{aj}(y).
\end{aligned}
\]
Indeed, for \(a=0\) the indicator selects \(z_0=y\) and the
\(z_1\)-sum equals one; for \(a=1\) it selects \(z_1=y\) and the
\(z_0\)-sum equals one.

The propensity weights reconstruct the treatment--covariate masses:
\[
p_j(P)\beta^A_j(a)=s_{aj}(P).
\]
For \(p_j(P)=0\), nonnegativity and
\(s_{0j}(P)+s_{1j}(P)=0\) give both arm masses zero. For
\(p_j(P)>0\), the two identities are
\[
p_j(P)e_j(P)=s_{1j}(P),\qquad
p_j(P)\{1-e_j(P)\}=p_j(P)-s_{1j}(P)=s_{0j}(P).
\]
Likewise,
\[
s_{aj}(P)\beta^Y_{aj}(y)=P(X=j,A=a,Y=y).
\]
If \(s_{aj}(P)=0\), both observed outcome atoms vanish. If
\(s_{aj}(P)>0\), the identities for success and failure are
\begingroup\small
\[
s_{aj}(P)\mu_{aj}(P)=q_{aj}(P),\qquad
s_{aj}(P)\{1-\mu_{aj}(P)\}
=s_{aj}(P)-q_{aj}(P)
=P(X=j,A=a,Y=0).
\]
\endgroup
Combining these calculations proves
\[
\Pr_{H_{\mathrm{ind}}}(X=j,A=a,Y=y)
=P(X=j,A=a,Y=y).
\]
Equality at every atom proves equality of the observed marginal with
\(P\), and in particular
\[
\Pr_{H_{\mathrm{ind}}}(X=j)=p_j(P).
\]
% lean: CausalSmith.Stat.AnnotationRarearmFrontier.independentExtension_observedMarginal

\item
The defining indicator gives
\[
\Pr_{H_{\mathrm{ind}}}
 (X=j,A=a,Y=y,Y_0=z_0,Y_1=z_1)=0
\qquad\text{whenever }y\ne z_a.
\]
Consequently,
\[
Y=AY_1+(1-A)Y_0
\qquad H_{\mathrm{ind}}\text{-almost surely},
\]
which is \cref{obj:ass:consistency}.
% lean: CausalSmith.Stat.AnnotationRarearmFrontier.independentExtension_consistency

\item
Summing over the factual outcome, and then using the unit sums of the
Bernoulli factors, gives
\[
\begin{aligned}
&\Pr_{H_{\mathrm{ind}}}(X=j,A=a,Y_0=z_0,Y_1=z_1)
=p_j(P)\beta^A_j(a)\beta^Y_{0j}(z_0)\beta^Y_{1j}(z_1),\\
&\Pr_{H_{\mathrm{ind}}}(X=j,A=a)=p_j(P)\beta^A_j(a),\\
&\Pr_{H_{\mathrm{ind}}}(X=j,Y_0=z_0,Y_1=z_1)
=p_j(P)\beta^Y_{0j}(z_0)\beta^Y_{1j}(z_1).
\end{aligned}
\]
Therefore, for every cell,
\[
\begin{aligned}
&\Pr_{H_{\mathrm{ind}}}(X=j,A=a,Y_0=z_0,Y_1=z_1)\,p_j(P)\\
&\quad=
\Pr_{H_{\mathrm{ind}}}(X=j,A=a)\,
\Pr_{H_{\mathrm{ind}}}(X=j,Y_0=z_0,Y_1=z_1).
\end{aligned}
\]
On an occupied cell, division by \(p_j(P)^2\) proves conditional
independence of \(A\) and \((Y_0,Y_1)\). More explicitly,
\[
\Pr_{H_{\mathrm{ind}}}
 (A=a,Y_0=z_0,Y_1=z_1\mid X=j)
=\beta^A_j(a)\beta^Y_{0j}(z_0)\beta^Y_{1j}(z_1).
\]
This is the stated product of three Bernoulli laws. Null cells have
zero probability, so the calculation establishes
\cref{obj:ass:exchangeability}. Together with the preceding steps,
\(H_{\mathrm{ind}}\) is a compatible extension.
% lean: CausalSmith.Stat.AnnotationRarearmFrontier.independentExtension_exchangeability

\item
Now fix any compatible extension \(H\). Define its potential-outcome
atom masses and its cellwise potential-success masses by
\[
\begin{aligned}
\rho^H_j(a,z_0,z_1)
&=\Pr_H(X=j,A=a,Y_0=z_0,Y_1=z_1)\\
&=\sum_y\Pr_H(X=j,A=a,Y=y,Y_0=z_0,Y_1=z_1),\\
\eta^H_{aj}
&=\sum_{a',z_0,z_1}\rho^H_j(a',z_0,z_1)\,z_a.
\end{aligned}
\]
These definitions apply to every cell and both arms. Expanding the
finite expectation, summing out \(Y\), and distributing the
difference gives
\[
\begin{aligned}
\mathbb E_H[Y_1-Y_0]
&=\sum_{j=1}^{d}\sum_{a,y,z_0,z_1}
 \Pr_H(X=j,A=a,Y=y,Y_0=z_0,Y_1=z_1)(z_1-z_0)\\
&=\sum_{j=1}^{d}\bigl(\eta^H_{1j}-\eta^H_{0j}\bigr).
\end{aligned}
\]
We identify the two success masses in each summand.
% lean: CausalSmith.Stat.AnnotationRarearmFrontier.poContrast_eq_sum_potentialSuccessMass

\item
Since the observed marginal of \(H\) equals \(P\), summing its atom
masses gives
\[
p_j(P)=\sum_{a',z_0,z_1}\rho^H_j(a',z_0,z_1),
\qquad
s_{aj}(P)=\sum_{z_0,z_1}\rho^H_j(a,z_0,z_1).
\]
By \cref{obj:ass:consistency},
\[
\Pr_H(X=j,A=a,Y=1-z_a,Y_0=z_0,Y_1=z_1)=0.
\]
Thus summing the factual-success atoms selects precisely those
potential-outcome pairs with \(z_a=1\), yielding
\[
q_{aj}(P)=\sum_{z_0,z_1}\rho^H_j(a,z_0,z_1)\,z_a.
\]
Conditional exchangeability in
\cref{obj:ass:exchangeability} supplies the atom identity
\[
\rho^H_j(a,z_0,z_1)\,p_j(P)
=
s_{aj}(P)\sum_{a'}\rho^H_j(a',z_0,z_1).
\]
For \(p_j(P)>0\), this follows by multiplying the conditional
factorization by \(p_j(P)^2\). For \(p_j(P)=0\), all the corresponding
atom masses are zero, so the same identity holds.

Multiplying this identity by \(z_a\) and summing over the potential
outcomes proves
\[
\begin{aligned}
q_{aj}(P)p_j(P)
&=\sum_{z_0,z_1}
 \rho^H_j(a,z_0,z_1)p_j(P)\,z_a\\
&=s_{aj}(P)\sum_{z_0,z_1}\sum_{a'}
 \rho^H_j(a',z_0,z_1)\,z_a\\
&=s_{aj}(P)\eta^H_{aj}.
\end{aligned}
\]
% lean: CausalSmith.Stat.AnnotationRarearmFrontier.compatible_success_identity

\item
Fix \(j\) and \(a\). First suppose \(p_j(P)=0\).
Nonnegativity of the atom masses and \(0\le z_a\le1\) give
\[
0\le\eta^H_{aj}
\le\sum_{a',z_0,z_1}\rho^H_j(a',z_0,z_1)
=p_j(P)=0.
\]
Hence
\[
\eta^H_{aj}=0=p_j(P)\mu_{aj}(P).
\]

Next suppose \(p_j(P)>0\). Model membership in
\cref{obj:def:model} supplies \cref{obj:ass:overlap}, so
\[
s_{1j}(P)=p_j(P)e_j(P)\ge\epsilon p_j(P)>0,
\qquad
s_{0j}(P)=p_j(P)\{1-e_j(P)\}\ge\epsilon p_j(P)>0.
\]
Dividing the preceding success identity by the positive arm mass
therefore gives
\[
\eta^H_{aj}
=p_j(P)\frac{q_{aj}(P)}{s_{aj}(P)}
=p_j(P)\mu_{aj}(P).
\]
The two cases establish this identity for every cell and both arms.
% lean: CausalSmith.Stat.AnnotationRarearmFrontier.compatible_cell_mean

\item
Substituting first the treatment success mass and then the control
success mass into the contrast expansion yields
\[
\begin{aligned}
\mathbb E_H[Y_1-Y_0]
&=\sum_{j=1}^{d}
 \bigl\{p_j(P)\mu_{1j}(P)-p_j(P)\mu_{0j}(P)\bigr\}\\
&=\sum_{j=1}^{d}p_j(P)
 \bigl\{\mu_{1j}(P)-\mu_{0j}(P)\bigr\}
=\tau(P),
\end{aligned}
\]
where the last equality is \cref{obj:def:target}. Since \(H\) was
arbitrary, this identifies the contrast for every compatible
extension.
% lean: CausalSmith.Stat.AnnotationRarearmFrontier.compatible_poContrast
\end{enumerate}
\end{proof}
